\documentclass{article}
\usepackage[margin=1in]{geometry}
\usepackage{graphicx}
\usepackage{setspace}
\setlength{\parindent}{0pt}
\setlength{\parskip}{1em}
\usepackage{amsmath}
\usepackage{algorithm}
\usepackage{algpseudocode}

\usepackage[backend=biber,style=numeric]{biblatex}
\addbibresource{Ex1.bib}

\title{A Photon Kinetic Framework for the OSIRIS Particle-in-Cell Code}
\author{Leah Hartman}

\begin{document}

\maketitle

Light propagating through density gradients of a plasma wakefield can undergo
significant frequency upshifting through a process known as photon acceleration.
The achievable frequency shift in this scheme is fundamentally limited only by
the depletion of the drive beam, with previous 1D simulations~\cite{sandberg2023prl} showing that
sufficiently energetic particle beam drivers can sustain the wake over distances
long enough to produce shifts exceeding 40 times, accompanied by an increase in 
pulse energy; quasi-3D simulations of the same scheme~\cite{sandberg2024pre} demonstrated shifts up
to 10 times, limited by simulation resolution and driver evolution rather than
the underlying physics. Capturing this effect in particle-in-cell (PIC) 
simulations is computationally demanding, as resolving the laser's electromagnetic
oscillations directly imposes strict resolution requirements via the 
Courant-Friedrichs-Lewy (CFL) condition.

We address this by implementing a photon kinetic model within the PIC code OSIRIS, 
in which the laser pulse is instead represented as a distribution of macro-particles,
each characterized by a position $\mathbf{x}(t)$ and wave vector $\mathbf{k}(t)$. The 
equations of motion follow directly from Hamilton's equations associated with the
Hamiltonian $\omega(\mathbf{x},\mathbf{k}, t)$~\cite{mendonca2001},
\begin{subequations}
\label{eq:hamilton}
\begin{align}
    \frac{d\mathbf{x}}{dt} &= \frac{\partial \omega}{\partial \mathbf{k}}, \label{eq:hamx}\\
    \frac{d\mathbf{k}}{dt} &= -\frac{\partial \omega}{\partial \mathbf{x}}. \label{eq:hamk}
\end{align}
\end{subequations}

For electromagnetic wave propagation in a cold, unmagnetized plasma, the local
dispersion relation may be written as
\begin{equation}
    \omega(\mathbf{x},\mathbf{k}, t) = \sqrt{|\mathbf{k}|^2 + \omega_p^2(\mathbf{x},t)}.
    \label{eq:dispersion}
\end{equation}
Substituting~\eqref{eq:dispersion} into~\eqref{eq:hamx} and~\eqref{eq:hamk} gives the
group velocity and wave-vector evolution used in the particle push,
\begin{equation}
    \frac{\partial \omega}{\partial \mathbf{k}} = \frac{\mathbf{k}}{\sqrt{|\mathbf{k}|^2 + \omega_p^2(\mathbf{x}, t)}},
\end{equation}
while the evolution of the wave vector is governed by spatial variations in
$\omega_p^2(\mathbf{x},t)$ and takes the form
\begin{equation}
    -\frac{\partial \omega}{\partial \mathbf{x}}
    =
    -\frac{1}{2\sqrt{|\mathbf{k}|^2 + \omega_p^2(\mathbf{x},t)}} \nabla \omega_p^2(\mathbf{x},t).
\end{equation}

These coupled equations of motion are advanced in time using a leapfrog-style
scheme, in which $\mathbf{x}$ and $\mathbf{k}$ are updated on staggered
half-steps, as outlined in Algorithm~\ref{alg:push}.

\newpage
\begin{algorithm}
\caption{Leapfrog update of a photon macro-particle's position $\mathbf{x}$ and 
wave vector $\mathbf{k}$ at each timestep, using the local plasma frequency 
$\omega_p^2$ and its gradient gathered from the grid.}
\label{alg:push}
\begin{algorithmic}[1]
\Require $\{\mathbf{x}_0^{(p)}, \mathbf{k}_0^{(p)}\}_{p=1}^{N}$ for $N$ macro-particles; grid field $\omega_p^2(\mathbf{x}, t)$
\Ensure Updated $\{\mathbf{x}_n^{(p)}, \mathbf{k}_n^{(p)}\}$ at each timestep $n$
\For{each timestep $n$}
    \For{each macro-particle $p = 1, \dots, N$}
        \State Gather $\omega_p^2(\mathbf{x}_n^{(p)}, t_n)$ and $\nabla \omega_p^2(\mathbf{x}_n^{(p)}, t_n)$ from the grid
        \State $\displaystyle \mathbf{k}_{n+1/2}^{(p)} \gets \mathbf{k}_{n-1/2}^{(p)}
            - \frac{\Delta t}{2\sqrt{|\mathbf{k}_{n-1/2}^{(p)}|^2 + \omega_p^2(\mathbf{x}_n^{(p)}, t_n)}}
            \,\nabla \omega_p^2(\mathbf{x}_n^{(p)}, t_n)$
        \State $\displaystyle \mathbf{x}_{n+1}^{(p)} \gets \mathbf{x}_n^{(p)}
            + \Delta t \, \frac{\mathbf{k}_{n+1/2}^{(p)}}{\sqrt{|\mathbf{k}_{n+1/2}^{(p)}|^2 + \omega_p^2(\mathbf{x}_n^{(p)}, t_n)}}$
    \EndFor
\EndFor
\end{algorithmic}
\end{algorithm}

By replacing the electromagnetic solve with this reduced kinetic description, the
model substantially reduces computational cost while retaining the relevant physics
of photon acceleration. 


\printbibliography
\end{document}
